\section{TEKNOLOGI}

\subsection{Daftar Teknologi}

Berikut adalah daftar teknologi yang digunakan dalam pengembangan Savora, sesuai dengan basis
kode aplikasi:

\begin{table}[htbp]
\centering
\caption{Daftar Teknologi Savora}
\label{tab:daftar_teknologi}
\small
\begin{tabularx}{\textwidth}{|l|X|X|}
\hline
\textbf{Kategori} & \textbf{Teknologi} & \textbf{Versi} \\
\hline
Bahasa Pemrograman & Go, JavaScript & 1.26, ES2024 \\
\hline
Frontend Framework & Next.js (App Router) + React & 16.x, 19.x \\
\hline
UI \& Styling & Tailwind CSS, lucide-react, Recharts & 4.x, 1.x, 3.x \\
\hline
Backend Framework & Go + Fiber & 1.26, v2.52 \\
\hline
Database & PostgreSQL & 16.x \\
\hline
ORM & GORM & 1.31 \\
\hline
Payment Gateway & Midtrans (mode Sandbox) & Snap \\
\hline
Food Trust Index \& Food Score & Rule-based (JavaScript) & - \\
\hline
Keyword Classification & Rule-based (JavaScript) & - \\
\hline
API Design & RESTful API (JSON) & - \\
\hline
Version Control & Git, GitHub & - \\
\hline
Hosting (rencana) & Vercel (frontend), Cloud/Container (backend) & - \\
\hline
\end{tabularx}
\end{table}

\subsection{Alasan Pemilihan Teknologi}

\subsubsection{Next.js + React untuk Frontend}
Next.js (dengan \textit{App Router}) dan React dipilih sebagai \textit{frontend framework} karena:
\begin{itemize}[leftmargin=*]
    \item \textbf{Komponen Reusable} : Arsitektur berbasis komponen memudahkan pembuatan UI yang konsisten dan mudah dipelihara.
    \item \textbf{Rendering Fleksibel} : Next.js mendukung \textit{server} dan \textit{client rendering} untuk performa dan SEO yang baik.
    \item \textbf{Developer Experience} : \textit{Routing} berbasis folder, \textit{fast refresh}, dan integrasi Tailwind CSS mempercepat pengembangan antarmuka \textit{mobile-first}.
    \item \textbf{Visualisasi Data} : Recharts digunakan untuk grafik pada dashboard analitik UMKM.
\end{itemize}

\subsubsection{Go + Fiber untuk Backend}
Go dengan framework Fiber dipilih karena:
\begin{itemize}[leftmargin=*]
    \item \textbf{Performa \& Konkurensi} : Go dikompilasi menjadi biner \textit{native} dan menangani banyak \textit{request} secara efisien dengan \textit{goroutine}.
    \item \textbf{Fiber yang Ringan} : Fiber menyediakan API bergaya Express di atas \textit{fasthttp}, cocok untuk REST API yang responsif.
    \item \textbf{Tipe Statis} : Menekan kelas kesalahan tertentu saat kompilasi sehingga layanan lebih andal.
\end{itemize}

\subsubsection{PostgreSQL + GORM}
\begin{itemize}[leftmargin=*]
    \item \textbf{PostgreSQL} : Database relasional yang robust dan mendukung ACID untuk data transaksional (produk, pesanan, ulasan, iklan).
    \item \textbf{GORM} : ORM untuk Go yang menyederhanakan pemetaan model dan \textit{auto-migration} skema database.
\end{itemize}

\subsubsection{Food Trust Index dan Food Score Berbasis Aturan}
Perhitungan Food Trust Index dan Food Score diimplementasikan sebagai fungsi \textit{rule-based} (deterministik)
di sisi frontend, sehingga skor dapat dihitung dan diperbarui \textit{real-time} mengikuti sisa waktu
rescue tanpa memerlukan layanan \textit{machine learning} terpisah. Rincian rumus dibahas pada bab
Algoritma dan Logika Aplikasi. Pendekatan berbasis \textit{machine learning} (mis. prediksi permintaan
dan \textit{dynamic pricing} otomatis) diposisikan sebagai \textit{roadmap} pengembangan lanjutan.

\subsubsection{Midtrans Sandbox untuk Pembayaran Cashless}
Savora menggunakan \textit{payment gateway} Midtrans pada mode \textit{sandbox} untuk mendemonstrasikan
alur pembayaran \textit{cashless} \cite{midtrans2025docs} tanpa transaksi uang asli. Pilihan ini
dipertimbangkan karena:
\begin{itemize}[leftmargin=*]
    \item \textbf{Realistis dan Dapat Diuji} : Mode \textit{sandbox} memungkinkan simulasi alur checkout,
    status pembayaran, dan \textit{webhook} secara \textit{end-to-end} untuk kebutuhan demo juri.
    \item \textbf{Model Pendapatan Jelas} : \textit{Service fee} 5\% otomatis ditambahkan per transaksi
    sebagai sumber pendapatan platform yang transparan.
    \item \textbf{Verifikasi Aman} : Status pembayaran hanya diperbarui setelah \textit{signature}/status
    dari notifikasi Midtrans diverifikasi backend, sehingga integritas transaksi terjaga.
\end{itemize}
Pengambilan makanan tetap dilakukan \textit{self-pickup} ke lokasi UMKM dengan validasi \textit{pickup
code} setelah pembayaran berstatus \textit{Paid}.

\subsection{Arsitektur Teknologi}

Savora menggunakan arsitektur \textit{client--server} tiga lapis: peramban pengguna menjalankan
aplikasi Next.js, yang berkomunikasi dengan backend Go/Fiber melalui REST API berformat JSON, dan
backend mengakses PostgreSQL melalui GORM. Gambar~\ref{fig:arsitektur} menunjukkan arsitektur
tingkat tinggi tersebut.

\begin{figure}[htbp]
\centering
\begin{tikzpicture}[
    font=\footnotesize,
    layer/.style={rounded corners, draw, thick, minimum width=3.1cm, minimum height=1cm, align=center},
    node distance=0.9cm and 1.2cm,
]
% Client layer
\node[layer, fill=blue!8] (browser) {Peramban Pengguna\\\scriptsize(Customer / UMKM / Admin)};
\node[layer, fill=blue!14, below=of browser] (front) {Frontend\\\textbf{Next.js 16 + React 19}\\\scriptsize Tailwind CSS, Recharts};
\node[layer, fill=green!12, below=of front] (api) {Backend API\\\textbf{Go 1.26 + Fiber v2}\\\scriptsize REST / JSON, \texttt{/api/*}};
\node[layer, fill=orange!14, below=of api] (orm) {ORM\\\textbf{GORM}};
\node[layer, fill=red!10, below=of orm] (db) {Database\\\textbf{PostgreSQL}};

% Food Score annotation (client-side rule-based)
\node[layer, fill=yellow!22, right=2.4cm of front, minimum width=3.4cm] (score) {Food Trust Index \& Food Score\\\scriptsize rule-based (client-side)\\\scriptsize meluruh terhadap waktu};

% External payment gateway
\node[layer, fill=purple!12, right=2.4cm of api, minimum width=3.4cm] (pay) {Midtrans Sandbox\\\scriptsize payment gateway\\\scriptsize service fee 5\%};

\draw[-{Latex}] (browser) -- (front);
\draw[-{Latex}] (front) -- node[right]{\scriptsize HTTPS} (api);
\draw[-{Latex}] (api) -- node[right]{\scriptsize query} (orm);
\draw[-{Latex}] (orm) -- (db);
\draw[-{Latex}, dashed] (front) -- (score);
\draw[{Latex}-{Latex}, dashed] (api) -- node[above, font=\scriptsize]{token/webhook} (pay);

\begin{scope}[on background layer]
    \node[draw, dashed, rounded corners, fit=(browser)(front)(score), inner sep=6pt, label={[font=\scriptsize]above:Client (Peramban)}] {};
    \node[draw, dashed, rounded corners, fit=(api)(orm)(db), inner sep=6pt, label={[font=\scriptsize]below:Server}] {};
\end{scope}
\end{tikzpicture}
\caption{Arsitektur High-Level Savora}
\label{fig:arsitektur}
\end{figure}

Alur komunikasi sistem:
\begin{enumerate}[leftmargin=*]
    \item \textbf{Client Layer} : Peramban pengguna menjalankan aplikasi Next.js/React. Food Score dihitung di sisi klien secara \textit{rule-based} dan diperbarui mengikuti sisa waktu rescue.
    \item \textbf{Application Layer} : \textit{Request} diteruskan ke backend Go/Fiber melalui REST API (\texttt{/api/*}) yang memproses \textit{business logic} produk, pesanan, analitik, dan iklan.
    \item \textbf{Payment Gateway} : Untuk pembayaran \textit{cashless}, backend membuat token transaksi ke Midtrans \textit{sandbox} dan memperbarui status order setelah memverifikasi notifikasi/\textit{webhook}.
    \item \textbf{Data Access Layer} : GORM memetakan model ke tabel dan menjalankan kueri ke database.
    \item \textbf{Data Layer} : PostgreSQL menyimpan data persisten (UMKM, produk, pesanan, ulasan, iklan, pembayaran).
\end{enumerate}
